\ta{Highlight number of papers and attendees/ unique authors. Mention it was virtual, still got good response. Panel with more engagement with audience.  Talk and Panel on education and communicating uncertainty. Panel on AI and decision-making.}

\ta{\underline{Comments from the workshop chair (Arvind)}: we would like to see some critical reflection/evaluation from the organizers of both past iterations of the workshop and the need/goals of continuing to offer it.}

\ta{\underline{New possible activities this year}: (1) Breakout sessions (potential tracks: (a) data modeling and transformation, (b) communication of uncertainty to scientific and general audience, (c) Evaluation of uncertainty visualization , (d) AI and uncertainty), (2) Pedagogial material (e.g., invited talks that discuss the "must know" tools and techniques (e.g., probability, confidence intervals, colormapping  of probabilities, spaghetti plots etc.) to enble uncertainty interpretation for general audience.) }

\ta{Notes from the previous workshop.} 

\ta{\underline{Keynote (Daniel Weiskopf)} - Uncertainty modeling, propagation, perception, congintion, and user interaction are all nontrivial problems. Uncertainty should be considered as a data in its own right. Many models assume data are independent, but considering correlation is super important. Creating models of uncertainty perception, cognition is much harder. How uncertainty propagation through nonlinear transformations can be tracked (through solutions other than Monte Carlo). Can AI be used for uncertainty visualization. How uncertainty visulization can be simplified for nonexpert users. Many existing solutions are application-specific. But we need more general solutions for broader impact. AI/ML advances provide great opportunity for generalizability of solutions, e.g., use AI to generate uncertainty models of input data/or use AI surrogates for uncertainty propagation. AI models themeselves are not perfect. Can uncertainty visualization of AI models be derived to support explainable AI?}

\ta{\underline{General notes from long/short talks} - (1) Tensor/high-dimensional/multivariate uncertainty vis all are challenging. (2)Surrogate models often used to understand the relationship between input-output simulation parameters. Uncertainty vis can help to steer simulation decisions, e.g., deciding which space of parameters needs more simulation. Communicating uncertainty of surrogate models is crucial. (3) User studies crucial to understand how uncertainty is interpreted by users. E.g., one study showed providing narration to describe the uncertainty can lead to poor/suboptimal decisions. Other study showed showing more samples of uncertainty in forecast can lead to enhanced decisions, but users are not willing to pay any extra costs associated to fetch forecast uncertainty info. (4) AI to accelerate uncertainty computations/forecasts (5) Complexity of visualization can significantly grow due to encoding of uncertainty, e.g., encoding uncertainty in glyphs can aggravate the problem of visual clutter. Howe to integrate uncertainty in vis with minimal cognitive overahead for users.}

\ta{\underline{General notes from lightening talks} - (1) Uncertainty of across visualization models difficult to portray/decode. (2) Uncertainty of neural networks used for vis. (3) Need new visual designs (glyphs) to effectively portray uncertainty. (4) Another open question is what is the best uncertainty model for visualization. E.g., For histograms, how many bins to use or is Gaussian best model? Evaluation of uncertainty models super important. (5) Interactive interfaces/toolboxes to enable users to explore uncertainty in the data (e.g., simulation of wildfire spread)}

\ta{\underline{General notes from the panel session} ( Han-Wei, Lace Padilla, Gerik Scheuermann, Matthew Kay): \textbf{(1) What is the main use of uncertainty? Is it trustworthiness, communication, discovery?}- (Han-Wei) Use of AI in science is limited because no one believes results. So main purpose is trust .  (Lace )Conveying trust may not be the main goal, but calibrating trust with accuracy of model should be the main goal?. How people imagine uncertainty in their mind is inconsistent across people. So displaying uncertainty anchors users to interpret uncertainty in a consistent manner. (Matt)Trying to help people make better decisions. As long as people make better and consistent decisions, that is what we ultimately want. (Gerik) Central point of uncertainty is avoid mistakes and help people making good decisions backed by the data ("being certain about uncertainty vis across users"). You want to understand significant uncertainties. You want to avoid decisions that are not backed by the data. \textbf{(2) Challenges in communication of uncertainty with scientific and general (non-experts) audience?} (Lace) Right type of vis is going to based on users, data, and task at hand. We need to make evaluation of if a vis supports users. Often user response to vis is unexpected. Generalizing fixed method across all users may not be possible as of now. (Matt) We don't have overarching model that says this uncertainty vis model is best for this situation beacause it is always complicated by background knowledge people have (e.g., cone of hurricane uncertainty can be misinterpreted as hurricane size. Uncertainty may be communicated with non-experts (e.g., news paper readers) by simplifying languat. This event will occur with 9/10 chance etc. and showing different scenarios. (Han-Wei) How can you bring humans in the decision-loop? You don't just display uncertainty to people, but allow users to interact/do something with it. Images may not be static. \textbf{(3) Chris: experience from geographic vis expert Alan M. MacEachren) Showing low/high uncertainty is good enough. If you try to convey more, then you can easily confuse people (have minimal descriptions). Do you agree or not?} (Gerik) we need little more than low/high. For example, showing detailed hurricane tracks can be useful. Uncertainty is certainly not making interpretation easier. What do you do with uncertainty is important. Depending on if its discovery or decision-making, actions can be different. (Lace) Just telling people high/low uncertainty can bias their decisions. Maybe showing high-quality sophisticated uncertainty vis may not be useful. Perhaps sending a text that say low/medium/high uncertainty about hurricane can help people to make evacuation decision? (Matt) expert in psychology may think people make terrible decions under uncertainty. In contrast, HCI experts think people do quite well under uncertainy. So uncertainty vis needs to be tailored to how people think, which can lead them to make better decision. \textbf{(4) Many times we think of uncertainty as probability, statics, which non-experts or even sometimes scientist not good at interpreting. How do we covercome that? Should we educate people about probability, statistics? Or should we change vis according to what people think (as per Matt)?}: (Han-Wei) People generally understand statistics, but they don't understand what 85\% uncertainty means. How can generative AI can help about how uncertain scenario might look. Present different scenarios to user as a means of easy decoding. \textbf{Most important area in uncertainty vis domain that needs to be tackled now:} (Gerik) How to process and communicate uncertainty and different scenarios. (Lace) Studying uncertainty and its connection to decision making in more generalized manner (outside clinical/lab settings etc.) (Matt) Connecting probabilistic programming with VIS. Communicating multiple levels of uncertainty simultaneously (e.g., probability, qualitative uncertainty)/ multiple scenarios for better decision making. (Han-Wei) How to visually encode uncertainty for ML/AI/data models. How to bring humans in loop for effective decision-making.}

\ta{\underline{Capstone (Jessica Hullman)} - Whats the value of better uncertainty visualization (e.g., whats the value of better representation of variance in exploratory data analysis, ML predictions etc.) Default approach is empirical study. Ask participants to make decisions under uncertainty to evaluated different uncertainty vis. Emprirical results, however, can be ambiguous. Another way is statistical theory (this should play a bigger role in uncertainty vis). Making general statements about uncertainty are difficult (e.g., gradient plots did well in our study and so should be used across similar task). If all uncertainty display options are good or if all are bad, it is difficult to choose one over another for generalizability. Bayesian theory can help to disambiguate whats the best/worst performance of uncertainty displays we could achieve to be able to evaluate them (Rational agents). Rational agents to derive best/worst case baselines. Can there be distribution-free uncertainty quantification approach?}